% Part: proof-theory
% Chapter: natural-deduction
% Section: introduction

\documentclass[../../../include/open-logic-section]{subfiles}

\begin{document}

\olfileid{pt}{nat}{int}

\olsection{Introduction}

Natural deduction is a family of !!{proof} systems in which each
logical connective or quantifier has a pair of inferences, one that
``introduces'' the operator (!!a{introduction} rule) and one that
``eliminates'' the operator (!!a{elimination} rule). For instance,
\Intro{\lif} has !!{formula}s of the form~$!A \lif !B$ as conclusion,
and \Elim{\lif} has !!{formula}s of the form~$!A \lif !B$ as premises.

The first two versions of natural deduction were introduced in the
1930s by Gerhard Gentzen and Stanisław Jaśkowski. Gentzen's is the
system proof theorists discuss mostly, and Jaśkowski gave rise to the
systems most widely used in teaching. In both systems, one can start
with \emph{assumptions}---!!{formula}s that don't need to be
!!{prove}d, but that can be used in !!{prove}ing other !!{formula}s.
An entire !!{proof} may depend on such assumptions, and since those
are unproven, the !!{proof} does not establish its conclusion (its
\emph{end-!!{formula}}), but only that its conclusion follows from the
assumptions.

Some inference rules are such that their conclusions no longer depend
on some assumptions: they ``discharge'' assumptions. For instance, the
\Intro{\lif} rule takes !!a{proof} of the conclusion~$!B$ from an
assumption~$!A$, and has $!A \lif !B$ as its conclusion. The !!{proof}
ending in $!A \lif !B$ no longer depends on~$!A$; the assumption $!A$
is discharged. In Gentzen's systems, !!{proof}s are trees of
!!{formula}s, the assumptions are the topmost !!{formula}s, and rules
like~\Intro{\lif} carry labels that indicate which assumptions are
discharged. In Jaśkowski's systems, parts of the !!{proof} that depend
on an assumption are set off in some way, e.g, indented behind a
vertical line or in a box, with the assumption~$!A$ discharged and the
consequent $!B$ of the conditional the first and last line in the box,
and the conclusion~$!A \lif !B$ of \Intro{\lif} outside the box.

The systems are ``natural'' because the inference rules mirror how we
(or mathematicians, at least) reason naturally. To establish a
hypothetical result (a conditional), we assume the antecedent and use
it to prove the consequent. That's~\Intro{\lif}. When we know a
conditional $!A \lif !B$ and its antecedent~$!A$, we conclude~$!B$.
That's~\Elim{\lif}. We use proof-by-cases to show that, assuming a
disjunction~$!A \lor !B$, some conclusion holds. That's~\Elim{\lor}.
To prove a general claim~$\lforall[x][!A(x)]$ we prove that $!A(c)$
holds, for an arbitrary~$c$. That's~\Intro{\lforall}. And so on.

For instance here's a simple !!{proof} of~$!A \lif (!A \lor !B)$ in
Gentzen-style natural deduction:
\[
\AxiomC{$\Discharge{!A}{x}$}
\RightLabel{\Intro{\lor}}
\UnaryInfC{$!A \lor !B$}
\DischargeRule{\Intro{\lif}}{x}
\UnaryInfC{$!A \lif (!A \lor !B)$}
\DisplayProof
\]
It starts with the assumption~$!A$, uses \Intro{\lor} to infer~$!A
\lor !B$, and then \Intro{\lif} to infer~$!A \lif (!A \lor !B)$. The
\Intro{\lif} inference discharges the assumption~$!A$; that's
indicated by the label~$x$.

Proof theorists prefer Gentzen's original systems working with trees
of !!{formula}s, although one variant is also important in proof
theory. !!^a{proof} of~$!A$ from assumptions~$\Gamma$ shows that $!A$
follows from~$\Gamma$, i.e., $\Gamma \Entails !A$. The inference rules
of natural deduction can naturally be read as telling us something
about such entailments, e.g., \Intro{\lif} says that if $\Gamma + !A
\Entails !B$ then $\Gamma \Entails !B$. It is only natural to
formulate the rules directly so that they operate on both the
assumptions and the conclusion that follow from them in the premises
and conclusion of the inference rule, i.e., they operate on sequents
$\Gamma \Sequent !A$. The simple proof above would look like this in
such a system:
\[
\Axiom$x:!A \fCenter !A$
\RightLabel{\Intro{\lor}}
\UnaryInf$x:!A \fCenter !A \lor !B$
\DischargeRule{\Intro{\lif}}{x}
\UnaryInf$\fCenter !A \lif (!A \lor !B)$
\DisplayProof
\]
As you see, the rules are the same: they work on the !!{formula} to
the right of~$\Sequent$, the !!{formula}s on the left just list the
assumptions they depend on at each step.

The complete set of !!{introduction}/!!{elimination} rule pairs is not
complete for classical logic by itself. We need to add two additional
rules that have to do with~$\lfalse$. The first is the
``intuitionistic absurdity rule''~\FalseInt, or ``explosion,'' which
says that from $\lfalse$ we can infer anything. The second is the
classical principle of indirect proof~\FalseCl{} (or the ``classical
absurdity rule'') which says that if we can !!{prove}~$\lfalse$
from~$\lnot !A$, we have !!{prove}d~$!A$. If we leave out~\FalseCl, we
get a system appropriate for intuitionsitic logic. If we leave out
both, we have what's known as ``minimal logic.''

We will discuss Gentzen-style natural deduction for classical and
intuitionsitic logic. These are the systems \Log{N1c} and \Log{N1i}.
The systems \Log{N2c} and \Log{N2i} are the corresponding systems
working on sequents $\Gamma \Sequent !A$ rather than single
!!{formula}s~$!A$.

\end{document}
